\documentclass[11pt]{article}

\usepackage[T1]{fontenc}
\usepackage{lmodern}
\usepackage[letterpaper,margin=20mm]{geometry}
\usepackage{microtype}
\usepackage[hidelinks]{hyperref}

\input{metadata}

\hypersetup{
  pdftitle={Cover letter for \papertitle},
  pdfauthor={\paperauthorone},
  pdfsubject={Submission to the Journal of Financial Econometrics Special Issue}
}

\pagestyle{empty}
\setlength{\parindent}{0pt}
\setlength{\parskip}{5.8pt}
\setlength{\emergencystretch}{1em}
\hyphenpenalty=10000
\exhyphenpenalty=10000
\begin{document}

\fontsize{10.5pt}{12.35pt}\selectfont

\noindent\hfill 30 August 2026\par

\begingroup
\setlength{\parskip}{0pt}
\textbf{Editors and Guest Editors}

\textit{Journal of Financial Econometrics}

Special Issue: \textit{Machine Learning in Financial Econometrics}
\endgroup

\smallskip

Dear Editors and Guest Editors,

Please consider our manuscript, ``\textbf{\papertitle},'' for the
\textit{Journal of Financial Econometrics} Special Issue on \textit{Machine
Learning in Financial Econometrics}.

The paper addresses a central theme of the call: how modern machine-learning
representations can enter financial econometrics without sacrificing economic
interpretation.
Rather than using embeddings only as predictive features, we represent each
firm's information as a distribution and use target-anchored Wasserstein
barycentric reconstruction to construct a predetermined, directed interaction field.
That field enters a spatial factor model in which feedback receives an
explicit exposure-adjustment interpretation.

The framework builds on well-understood spatial-econometric and spatial
asset-pricing foundations, but moves one step upstream: instead of taking the
interaction matrix as given, we construct it from firms' semantic information
footprints.
To our knowledge, this offers a novel generalization of spatial financial
econometrics from geographic or supplied-network spaces to
distribution-valued semantic spaces while retaining an admissible and
economically interpretable interaction operator.

Empirically, all text-derived fields are constructed from 2018--2022
information and frozen before estimation on 2023--2026 returns.
The Wasserstein-barycentric field delivers stronger conditional fit than both
equal weighting of the same peer support and a conventional distance-decay
transformation of the same Wasserstein distances.
When the barycentric and persistent news co-mention fields enter jointly,
both remain incrementally informative.
We also evaluate the construction across a broad panel of modern open-weight
embedding models, dimensions, capacities, architectures, and information vintages.
These models are widely used in contemporary information retrieval and
retrieval-augmented-generation pipelines, making the robustness exercise
especially relevant to current ML practice.

We believe this combination of modern ML representation, interpretable
spatial financial econometrics, and reproducible formal verification is
particularly well aligned with the aims of the Special Issue.

Thank you for considering our manuscript.

\begingroup
\setlength{\parskip}{0pt}
Sincerely,\par
\smallskip
\paperauthorone\par
Department of Finance and Tax, University of Cape Town\par
\href{mailto:\paperemail}{\paperemail}
\endgroup

\end{document}
